\documentclass[12pt]{article}
\usepackage[utf8]{inputenc}
\usepackage[T1]{fontenc}
\usepackage{hyperref}
\title{Awu Language Resources for K-12 Curriculum Development}
\author{Compiled by Claude Sonnet 5.5 (claude-sonnet-5-5) for Andr\'as Kornai}
\date{}
\begin{document}
\maketitle

\section{Introduction}
Awu (ISO 639-3 yiu; Glottocode awuu1235), also called Lope or Luowu, is a Loloish language of Yunnan, China, closely related to Nisu. Wikipedia gives about 20,000 speakers as of 2002. Only overview and audio-directory material turned up in the searches.

\section{Primary Educational Resources}
No primers, textbooks or school programmes for Awu were found.

\section{Linguistic and Reference Materials}
\begin{itemize}
\item Wikipedia, Lope language: \url{https://en.wikipedia.org/wiki/Lope_language}
\item Wikipedia, list of lesser-known Loloish languages: \url{https://en.wikipedia.org/wiki/List_of_lesser-known_Loloish_languages}
\end{itemize}
No grammar or dictionary was found.

\section{Religious and Community Resources}
\begin{itemize}
\item Asia Harvest profile, Northern Awu: \url{https://people-groups.asiaharvest.org/people-group-profiles/yunnan/awu-northern.pdf}
\item Asia Harvest profile, Southeastern Awu: \url{https://people-groups.asiaharvest.org/people-group-profiles/yunnan/awu-southeastern.pdf}
\end{itemize}

\section{Digital and Audio Resources}
\begin{itemize}
\item Global Recordings Network, Awu: \url{https://globalrecordings.net/en/language/yiu}
\end{itemize}
This page lists audio and video material, which the search summary described as Bible stories and evangelistic messages. A Global Recordings page for ``Yi: Xiangtang'' is also listed in the results, but I did not verify that it is the same language: \url{https://globalrecordings.net/en/language/6254}

\section{Curriculum Development Context}
The language is used by all adults in the community but not by all young people, and it is not known to be taught in schools. Little exists beyond overview and audio material, so curriculum development would need new documentation work.
\end{document}